\subsubsection{Concept 3}

This design would eliminate cameras entirely from the sensing design and use a simple time-of-flight
distance sensor mounted on a rotating servo, allowing it to sweep across a line above the drone to detect a perch.

\paragraph{Budget}

\begin{table}[ht]
    \centering
    \caption{Estimated budget for the lidar sensing design.}
    \label{tab:lidar}


    \begin{spreadtab}{{tabular}{|c|c|c|c||c|}}
        \hline
        @Item                       & @Supplier                                                                                                         & @Unit Cost (USD) & @Quantity & @Total Cost (USD) \\
        \hline
        @SMA Wire                   & @\href{https://www.kelloggsresearchlabs.com/product/round-wire/}{Kelloggs Research Labs}                          & 15.99 & 1 & [-1,0] * [-2,0] tag(start) \\
        @Mechanical Materials and Fabrication & @-                                                                                                      &     0 & 1 & [-1,0] * [-2,0] tag(pretax)\\
        @Tax                        & @- & sum(cell(start):cell(pretax)) * 0.07 & 1 & [-1,0] * [-2,0] tag(posttax) \\
        @Misc. Spares               & @- & sum(cell(start):cell(posttax)) * 0.2 & 1 & [-1,0] * [-2,0] tag(end) \\
        \hline
        \multicolumn{4}{|c||}{@Total} & sum(cell(start):cell(end)) \\
        \hline
    \end{spreadtab}
\end{table}

\paragraph{Evaluation}

This method would be computationally simpler than the camera based systems, though more mechanically
complex and possibly heavier. Additionally it would be unable to detect the orientation of a cylindrical 
perch, due to only measuring a cross section of the object.
